\subsection{Cinemática diferencial e Jacobiano do efetuador}

Para relacionar as velocidades das juntas do manipulador robótico com a velocidade espacial\footnote{\,Velocidade dentro do espaço de trabalho do manipulador robótico} do efetuador, utiliza-se o Jacobiano $J(q)$; o qual mapeia as velocidades em cada junta e as transforma em velocidades - lineares e angulares - do efetuador. Considerando
$\vec q$ como o vetor de juntas e
$\dot{q}$ as suas respectivas velocidades, a velocidade espacial do efetuador (linear e angular) no frame da base é dada por:

%% ===== EXPLICAÇÃO =====
% O vetor do lado esquerdo (x_ponto) é a velocidade espacial do efetuador
% V_E é a velocidade linear (m/s)
% Omega_E é a velocidade angular (rad/s)

% O vetor q_ponto contém as velocidades das juntas.
% Cada elemento é thata_ponto_i para as juntas revolutas e d_ponto_i para as juntas prismáticas.

% O Jacobiano J(q) é a mtriz 6xn que mapeia velocidades de juntas para velocidades do efetuador.

% Ou seja: sabendo as velocidades das juntas, multiplica-se por J(q) e obtém como o efetuador está se movendo no espaço.

\begin{equation}
\label{eq:jac_map_vel1}
\dot{\vec{x}}_{E}=
\begin{bmatrix}
    \vec{v}_{E} \\
    \boldsymbol{\omega}_{E}
\end{bmatrix}
= J(q)\dot{q}
\end{equation}

Onde:
\begin{equation}
    \label{eq:jac_map_vel2}
J(q) \in \mathbb{R}^{6 \times n}
\end{equation}

%% ===== EXPLICAÇÃO =====
% Para cada junta "i", a coluna correspondente "J_i" é construida a partir do eixo da junta (z_i-1) e das posições de referencia (p_e, p_i-1).
% p_E = posição do efetuador
% p_i-1 = posição da junta anterior
% z_i-1 = eixo da junta anterior (eixo de rotação para junta revoluta, direção de movimento para junta prismática)
%
% Junta revoluta: o produto vetorial mostra que a velocidade linear do efetuador vem da rotação em torno de z_i-1; a velocidade angular é simplesmente o eixo da rotação.
%
% Junta prismática: não há rotação, portanto omega = 0; o efetuador se move linearmente na direção do eixo da junta.
%
% Portanto, cada coluna do Jacobiano corresponde à contribuição individual de uma junta para a velocidade do efetuador como um todo.
As colunas de $J(q)$ para determinada junta revoluta ($J_{i}$), (usando $z_{i-1}$ e o ponto $p_{i-1}$ obtidos das matrizes homogêneas, ${}^{0}\,T_{i-1}$):
\begin{equation}
    \label{eq:jac_columns_revoluta}
\text{juntas revolutas:}
\begin{cases}
J_{v,i} = z_{i-1} \times \big(p_{E} - p_{i-1}\big),\\
J_{\omega,i} = z_{i-1},
\end{cases}
\end{equation}

E para a junta prismática:
\begin{equation}
    \label{eq:jac_columns_prismatica}
\text{junta prismática:}
\begin{cases}
J_{v,i} = z_{i-1},\\
J_{\omega,i} = {0}.
\end{cases}
\end{equation}

%% ===== EXPLICAÇÃO =====
% Qdo o manipulador está preso a uma base não-fixa, a veloc. do efetuador não depende apenas das juntas: há uma contribuição da movimentação da base.

% O termo com a adjunta (Adj) leva em conta a contribuição do mov da base
% Se a base é fixa, omega = 0, sobra apenas o termo J(q)q_ponto.

% A Adj possui os termos: Rotação (R), posição (p) e o produto vetorial (p_x) em forma matricial, matriz anti-simetrica.

% Portanto, esse termo "transporta" a velocidade da base para o frame do efetuador.
Durante a fase operacional, quando o subsistema do manipulador robótico atua como um corpo não-rígido, a torção total no $frame \{0\}$ é resultado da contribuição da base e das juntas via a adjunta $Adj_{{}^{0}T_{B}}$:
\begin{equation}
    \label{eq:jac_com_base1}
{}^{0} \, \vec{V}_{E} =
Adj_{\,{}^{0}T_B}
\begin{bmatrix}
{0} \\
\vec{\omega}_b
\end{bmatrix}
+ J(q)\dot{q},
\end{equation}

Onde:
\begin{equation}
    \label{eq:jac_com_base2}
Adj_T =
\begin{bmatrix}
\mathbf{R} & 0\\
[\vec{p}]_\times \mathbf{R} & \mathbf{R}
\end{bmatrix}.
\end{equation}

%% ===== EXPLICAÇÃO =====
% Forma dual do Jacobiano.
% Se o efetuador aplica uima força (f_E) e um momento (m_E) sobre o ambiente, o Jacobiano transposto mapeia isso para os torques ou forças necessários em cada junta.
%
% Portanto, J leva as velocidades de juntas (veloc do efetuador)
% J^T leva forças no efetuador (torques nas juntas)
O mapeamento das forças e dos torques é dado por
\begin{equation}
    \label{eq:forca_para_torque1}
\boldsymbol{\tau} = J(q)^{T} \vec{w}_{E}
\end{equation}

Onde:
\begin{equation}
    \label{eq:forca_para_torque2}
\vec{w}_{E} =
\begin{bmatrix}
\vec{f}_{E} \\
\vec{m}_{E}
\end{bmatrix}.
\end{equation}

%% ===== EXPLICAÇÃO =====
% "m" é a dimensão da tarefga, 6 se for controle completo de posição e orientação.
% Se "rank(J)" cai, o robô perde graus de liberdade efetivos.
% >>> Exemplo: um braço esticado pode perder a capacidade de se mover em certas direções.
% Os comandos podem "explodir" proximo de singularidades.

% O teste numerico é um jeito de estimar numericamente o Jacobiano por diferenças finitas, desloca-se um pouco "q" e ve como x_e varia.
% Isso é util pra detectar singularidades.
\subsection{Singularidades}
A configuração $q$ é singular quando $J(q) < m$, onde $m$ é a dimensão da tarefa (por exemplo, $m=6$ para controle de torção completa). Uma checagem numérica simples pode usar diferenças finitas:
\begin{equation}
    \label{eq:jac_singularidadeCheck}
J(q) \approx
\frac{\vec{x}_E(q+\Delta q) - \vec{x}_E(q)}
{\Delta q}.
\end{equation}

%% RESUMO GERAL
% Jacobiano = derivada da cinemática direta.
% >>> Conecta velocidades de juntas com velocidades do efetuador.
% Colunas = efeito de cada junta.
% >>> Revoluta: rotação gera veloc linear (via produto vetorial)
% >>> Prismática: só deslocamento linear
% Base móvel = se o manipulador não está fixo, há termos adicionais via adjunta.
% Dualidade
% >>> J(q_ponto): velocidade do efetuador
% >>> J^T(w_e): torques ou juntas necessários
% Singularidades = O Jacobiano perde "posto" e o robô não consegue se mover em todas as direções desejadas.
